\origpage{383}{N. Robertson, P.D. Seymour / Journal of Combinatorial Theory, Series B 90 (2004) 325--385}
Thus there is no $e \in E(\Gamma_0)$ bordering $\bar{r}\setminus r$, not a short side, with $\beta(e) \subseteq ins(F)$, and so
\begin{equation*}
N(r^*) \cap ins(F) = V(\Phi_\chi) \cap ins(F).
\end{equation*}
But then the second outcome of the theorem holds, a contradiction. This proves (4).

Let $\mathscr{T}'$ be the set of all separations $(A', B')$ of $sk(\Gamma)$ of order $\le n$ such that there exists $(A, B) \in \mathscr{T}$ with $V(A) = V(A')$ and $V(B) = V(B')$. By Theorem 14.1 of [9], $\mathscr{T}'$ is a tangle in $sk(\Gamma)$ of order $n+1$ and $rep(\mathscr{T}') = n+1$, since $ord(\mathscr{T}) > \frac{3}{2}n$ and $rep(\mathscr{T}) \ge n+1$. \fix{Theorem 14.1 of [9] gives only that $\mathscr{T}'$ is a tangle; Theorem 3.2 of [7] needs also that $\mathscr{T}'$ is respectful, and we prove this now. If $\hat{\Sigma}_\chi$ is a sphere, every tangle in a drawing in $\hat{\Sigma}_\chi$ is respectful; so we assume it is not. Call an open set $W \subseteq \hat{\Sigma}_\chi$ \textit{flat} if any two paths in $W$ with the same ends are homotopic in $\hat{\Sigma}_\chi$ with their ends fixed. Let $F$ be an $sk(\Gamma)$-normal $O$-arc with $|F \cap V(\Gamma)| \le n$. As in the proof of (6.4) of [6], the union of the regions of $\Gamma$ in $\hat{\Sigma}_\chi$ and of small open discs around the vertices in $F \cap V(\Gamma)$ is flat; this uses that $|F \cap V(\Gamma)| \le n < rep(\mathscr{T})$, and that every circuit $C$ of $\Gamma^*$ with $\frac{1}{2}|E(C)| < rep(\mathscr{T})$ bounds a disc, namely $ins(C)$. The regions of $sk(\Gamma)$ are the regions of $\Gamma$ and the edges $c$ of $\Gamma$ with $|\tilde{c}| = 3$, each of which is an open disc. Adding these discs keeps the set flat: any two ends $a, b$ of such an edge $c$ lie in the boundary of a region $r$ of $\Gamma$ that has a part of $bd(\bar{c})$ between $a$ and $b$ in its boundary, so a path in $\bar{c}$ from $a$ to $b$ is homotopic to a path in $r \cup \{a, b\}$. Since $F$ is $sk(\Gamma)$-normal, it lies in this flat set, so $F$ is null-homotopic, and hence $F$ bounds a closed disc $\Delta$. Let $\Delta'$ be the closure of $\hat{\Sigma}_\chi \setminus \Delta$. Then $(sk(\Gamma) \cap \Delta, sk(\Gamma) \cap \Delta')$ is a separation of order at most $n$, so it or its reverse is in $\mathscr{T}'$, and we may assume that $(sk(\Gamma) \cap \Delta', sk(\Gamma) \cap \Delta) \in \mathscr{T}'$, for otherwise $\Delta$ is as required. So there exists $(A, B) \in \mathscr{T}$ with $V(A) = V(sk(\Gamma) \cap \Delta')$ and $V(B) = V(sk(\Gamma) \cap \Delta)$, of order at most $n < rep(\mathscr{T})$. As in the proof of (6.3) of [6], there is a subdrawing $N$ of $\Gamma^*$ such that every circuit $C$ of $N$ has $\frac{1}{2}|E(C)| < rep(\mathscr{T})$, and every edge of $A$ lies in $ins(C)$ for some circuit $C$ of $N$; these circuits are null-homotopic, so by (7.5) of [6] there is a closed disc $D_1$ including $U(N)$, and then $D_1$ includes each such $ins(C)$, since $\hat{\Sigma}_\chi$ is not a sphere. Hence every edge of $A$ lies in $D_1$. On the other hand, every $O$-arc in $U(sk(\Gamma) \cap \Delta)$ is null-homotopic, since it lies in the disc $\Delta$; so by the flat set above and the argument of claim (2) in the proof of (6.4) of [6] (which uses (7.5) of [6]), there is a non-null-homotopic $O$-arc $F' \subseteq U(sk(\Gamma) \cap \Delta') \setminus U(sk(\Gamma) \cap \Delta)$. Then $F' = U(C')$ for a circuit $C'$ of $sk(\Gamma)$. Each edge of $C'$ lies in an edge $c$ of $\Gamma$ which has an end not in $V(B)$, since $F'$ misses $U(sk(\Gamma) \cap \Delta)$; so $c$ is an edge of $A$, and $c \subseteq D_1$. Thus $F' \subseteq D_1$, and $F'$ is null-homotopic, a contradiction. So $\mathscr{T}'$ is respectful.}\fixnote{fix:15.2-respectful-proof-added}{The paper does not say why $\mathscr{T}'$ is respectful. The argument above is the one used in the formalization.}

\hspace*{12pt}(5) \textit{Let $r^*, N(r^*)$ be as in (4). There is no $(A', B') \in \mathscr{T}'$ with $N(r^*) \subseteq V(A')$ and $|N(r^*)| > |V(A' \cap B')|$.}

\textbf{Subproof.} Choose $(sk(\Gamma))^*$ such that $\Gamma^* \subseteq (sk(\Gamma))^*$. Suppose there is such an $(A', B')$. Then by 15.1, there is a circuit $F$ of $(sk(\Gamma))^*$ of length $\le 2n$, bounding a closed disc $\Delta$ with $r^* \in \Delta$, such that
\begin{equation*}
(sk(\Gamma) \cap \Delta, sk(\Gamma) \cap \Delta') \in \mathscr{T}',
\end{equation*}
where $\Delta'$ is the closure of $\hat{\Sigma}_\chi \setminus \Delta$, and
\begin{equation*}
|N(r^*) \cap \Delta| > \frac{1}{2}|E(F)|.
\end{equation*}
For every vertex $v$ of $F$ which lies in the interior of an edge of $\Gamma$, there is a vertex $v'$ of $\Gamma^*$ adjacent in $(sk(\Gamma))^*$ to both neighbours of $v$ in $F$; and by replacing $v$ by $v'$, and repeating, we may assume that $F \subseteq \Gamma^*$. But this contradicts (4). This proves (5).

\hspace*{12pt}(6) \textit{Let $e \in E(\Gamma_0)$ be internal. There is no disc $\Delta \subseteq \hat{\Sigma}_\chi$ with $bd(\Delta)$ $sk(\Gamma)$-normal, such that $|\tilde{\beta}(e) \cap \Delta| > |V(\Gamma) \cap bd(\Delta)|$ and $(\Gamma \cap \Delta, \Gamma \cap \Delta') \in \mathscr{T}'$, where $\Delta'$ is the closure of $\hat{\Sigma}_\chi \setminus \Delta$.}

\textbf{Subproof.} Suppose that $\Delta$ is such a disc. Then $|V(\Gamma) \cap bd(\Delta)| \le 2$, since $|\tilde{\beta}(e) \cap \Delta| \le 3$. Since $d(\beta(e), r^*) > n$ for every pole $r^*$, it follows that $\Delta \subseteq \Sigma_\chi \setminus bd(\Sigma_\chi)$. Now
\begin{equation*}
|V(\Gamma) \cap \Delta| \ge |\tilde{\beta}(e) \cap \Delta| > |V(\Gamma) \cap bd(\Delta)|
\end{equation*}
and so $\Delta \cap V(\Gamma) \nsubseteq bd(\Delta)$. Since $\Gamma$ is internally 3-connected and hence so is $sk(\Gamma)$, we deduce that $\Delta'$ is a disc and $\Delta' \cap V(\Gamma) \subseteq bd(\Delta')$. But then the third tangle axiom is violated, because $V(\Gamma \cap \Delta) = V(\Gamma)$. This proves (6).

Let $n$ be so large (in terms of $\hat{\Sigma}_\chi$ and $\Gamma_0$) that Theorem 3.2 of [7] applies. From (1), (3), (4), (6), and Theorem 3.2 of [7] applied to $\mathscr{T}'$ and $sk(\Gamma)$, the trees $T_v$ in (3) may be chosen to be subgraphs of $sk(\Gamma)$. For each $v \in V(\Gamma_0)$, let $\sigma(v)$ be the induced subgraph of $sk(\Gamma)$ with vertex set $V(T_v)$, and for each $e \in E(\Gamma_0)$ let $\sigma(e) = \beta(e)$. Then $\sigma$ is an inflation of $(\Gamma_0, \phi_0)$ in $(\Gamma, \phi)$, a contradiction, as required. $\square$
